\documentclass[11pt]{article}

\usepackage{amsmath, amssymb, amsthm}
\usepackage{geometry}
\usepackage{hyperref}
\usepackage{algorithm}
\usepackage{algpseudocode}
\usepackage{booktabs}
\usepackage{bm}

\geometry{margin=1.25in}

\newtheorem{remark}{Remark}
\newtheorem{proposition}{Proposition}

\title{%
  A Discretize-First Linearized ADMM Solver\\
  for the Dynamic Schr\"{o}dinger Bridge
}
\author{}
\date{}

\begin{document}

\maketitle
\tableofcontents
\bigskip

%============================================================
\section{Introduction}
%============================================================

This document describes the design and implementation of a 1D solver for the
\emph{dynamic Schr\"{o}dinger Bridge} (SB) problem using a two-grid
linearized Alternating Direction Method of Multipliers (ADMM).  Three
intertwined design choices shape the entire solver:
\begin{enumerate}
  \item \textbf{Discretize first, then optimize.}  The continuous PDE
        constraint is discretized on a grid before any optimality conditions
        are derived.  ADMM is applied to the resulting finite-dimensional
        problem.
  \item \textbf{Two staggered grids.}  The objective and the constraint live
        most naturally on different grids.  We maintain both explicitly and
        couple them via interpolation operators inside ADMM.
  \item \textbf{Implicit diffusion discretization.}  The diffusion term in
        the Fokker-Planck constraint is treated implicitly, rendering the
        scheme unconditionally stable and avoiding the parabolic CFL
        restriction $\Delta t = O(\Delta x^2)$.
\end{enumerate}
Each choice is elaborated below after introducing the continuous problem.

%============================================================
\section{The Continuous Problem}
%============================================================

\subsection{Benamou--Brenier Optimal Transport}

The \emph{Benamou--Brenier} (BB) formulation of optimal transport seeks a
density--momentum pair $(\rho, m)$ on the space-time cylinder
$(t,x) \in [0,1]^2$ that minimizes the kinetic energy
\begin{equation}
  J(\rho, m)
  = \int_0^1 \int_0^1 \frac{1}{2}\,\frac{m(t,x)^2}{\rho(t,x)}\,dx\,dt
  \label{eq:BB_objective}
\end{equation}
subject to the continuity equation with prescribed marginals,
\begin{equation}
  \partial_t \rho + \partial_x m = 0,
  \qquad
  \rho(0,\cdot) = \rho_0,\quad
  \rho(1,\cdot) = \rho_1,\quad
  \partial_x \rho\big|_{x=0,1} = 0.
  \label{eq:continuity}
\end{equation}
The integrand $m^2/(2\rho)$ is the kinetic energy density of a fluid with
density $\rho$ and momentum $m = \rho v$ where $v$ is the velocity field.

\subsection{The Schr\"{o}dinger Bridge}

Adding entropy regularization replaces the continuity equation with the
\emph{Fokker--Planck equation},
\begin{equation}
  \partial_t \rho + \partial_x m = \varepsilon\,\partial_{xx}\rho,
  \qquad
  \rho(0,\cdot) = \rho_0,\quad \rho(1,\cdot) = \rho_1,
  \label{eq:FP}
\end{equation}
with the same objective~\eqref{eq:BB_objective}.  The parameter
$\varepsilon \geq 0$ controls the strength of diffusion:
\begin{itemize}
  \item $\varepsilon = 0$: recovers the pure OT (Benamou--Brenier) problem.
  \item $\varepsilon > 0$: the Schr\"{o}dinger Bridge; paths are smoother
        and more diffuse, corresponding to the most likely trajectory of
        $N$ Brownian particles conditioned on their initial and final
        empirical distributions.
\end{itemize}
The term $m = \rho v$ retains its role as the momentum, but $v$ is now a
drift that guides the diffusive particles.

%============================================================
\section{Discretize First, Then Optimize}
%============================================================

\subsection{Two Philosophies}

A PDE-constrained optimization problem can be attacked in two orders:

\begin{description}
  \item[Optimize-then-discretize (OtD):]
    Derive the infinite-dimensional optimality conditions (Hamilton--Jacobi--Bellman
    equation, adjoint system, etc.) and only then discretize those PDEs for
    numerical solution.  The resulting scheme is often non-trivially coupled
    and requires careful adjoint consistency.

  \item[Discretize-then-optimize (DtO):]
    Directly discretize the state equation and the objective on a grid,
    producing a finite-dimensional constrained optimization problem.  An
    off-the-shelf algorithm (here ADMM) is then applied to this discrete
    problem.
\end{description}

We adopt the \textbf{DtO} approach.  The advantages in this context are:
\begin{itemize}
  \item ADMM's convergence guarantees apply to the finite-dimensional problem
        without needing to track consistency of an adjoint PDE discretization.
  \item The prox and projection sub-steps are posed on a fixed grid and can
        be computed exactly or near-exactly at each iteration.
  \item The implicit diffusion treatment (Section~\ref{sec:implicit}) emerges
        naturally from the grid formulation rather than requiring separate
        analysis at the continuous level.
\end{itemize}

\subsection{The Discrete Problem}

Let $N_t$ and $N_x$ be the numbers of time and space cells, with
$\Delta t = 1/N_t$ and $\Delta x = 1/N_x$.  Cell centers are at
$(t_{i+1/2}, x_{j+1/2})$ and cell edges (faces) are at integer nodes.

We discretize $(\rho, m)$ on a space-time grid and require the discrete
analogue of~\eqref{eq:FP} to hold at every grid point.  The objective is
replaced by
\begin{equation}
  J_h(\rho, m)
  = \Delta t\,\Delta x \sum_{i,j} \frac{m_{ij}^2}{2\,\rho_{ij}},
  \label{eq:discrete_obj}
\end{equation}
with $J_h = +\infty$ whenever any $\rho_{ij} \leq 0$.

%============================================================
\section{Two Grids and Why We Need Them}
%============================================================

\subsection{Motivation: Conflicting Requirements}

The objective~\eqref{eq:discrete_obj} requires $\rho$ and $m$ to live at
the \emph{same} spatial and temporal locations, so their ratio $m^2/\rho$
can be formed pointwise.

The Fokker--Planck equation, on the other hand, is a conservation law of the
form $\partial_t\rho + \partial_x m = \varepsilon\partial_{xx}\rho$.  A
\emph{conservative} finite-difference discretization naturally lives on a
staggered grid: the density $\rho$ at cell centers, the flux $m$ at cell
faces.  Using the same collocated grid for both $\rho$ and $m$ in the
constraint produces a checkerboard-mode instability and requires artificial
stabilization.

These two requirements are in tension:
\begin{itemize}
  \item \textbf{Objective}: $\rho$ and $m$ collocated at cell centers.
  \item \textbf{Constraint}: $\rho$ at integer time nodes $\times$ cell
        centers, $m$ at half-integer times $\times$ cell faces.
\end{itemize}

\subsection{Grid Definitions}

We therefore maintain \emph{two} sets of variables:

\paragraph{Grid 1 (collocated), variables $(\rho, m)$:}
Both live at space-time cell centers $(t_{i+1/2}, x_{j+1/2})$.
\begin{align*}
  \rho_{ij} &\approx \rho(t_{i+1/2},\, x_{j+1/2}), \quad
    \text{shape } N_t \times N_x,\\
  m_{ij}    &\approx m(t_{i+1/2},\, x_{j+1/2}), \quad
    \text{shape } N_t \times N_x.
\end{align*}

\paragraph{Grid 2 (staggered), variables $(\bar\rho, b)$:}
Density at integer times, flux at spatial faces.
\begin{align*}
  \bar\rho_{ij} &\approx \rho(t_i,\, x_{j+1/2}),\quad
    \text{shape } (N_t-1) \times N_x
    \quad\text{(interior time nodes; BCs absorbed into } d\text{)},\\
  b_{ij}        &\approx m(t_{i+1/2},\, x_j),\quad
    \text{shape } N_t \times (N_x - 1)
    \quad\text{(interior faces; Neumann BCs enforced)}.
\end{align*}
The boundary values $\bar\rho^0 = \rho_0$ and $\bar\rho^{N_t} = \rho_1$ are
not stored as unknowns; they appear on the right-hand side of the constraint.

\subsection{Interpolation Operators}

The two grids are linked by averaging operators:
\begin{align}
  (A_\rho\,\rho)_i     &= \tfrac{1}{2}(\rho_i + \rho_{i+1}), \quad
    i = 0,\ldots,N_t-2,
    \label{eq:Arho}\\
  (A_m\, m)_j          &= \tfrac{1}{2}(m_j + m_{j+1}), \quad
    j = 0,\ldots,N_x-2.
    \label{eq:Am}
\end{align}
Stacking these, the interpolation from Grid 1 to Grid 2 is written
$A(\rho,m) = (A_\rho\,\rho,\; A_m\,m)$.

\subsection{The ADMM Splitting}

ADMM splits the problem by introducing duplicates:
\begin{equation}
  \min_{(\rho,m),\,(\bar\rho,b)}
  J_h(\rho, m) + \iota_C(\bar\rho, b),
  \qquad
  \text{subject to}\quad
  A(\rho, m) = (\bar\rho, b),
  \label{eq:split}
\end{equation}
where $\iota_C$ is the indicator of the Fokker--Planck constraint set
\begin{equation}
  C = \bigl\{(\bar\rho, b) \mid
    (I_x \otimes D_t - \varepsilon L_x \otimes A_t)\bar\rho
    + (D_x \otimes I_t)\,b = d\bigr\}.
  \label{eq:C}
\end{equation}
Each sub-problem is now tractable:
\begin{itemize}
  \item The $(\rho, m)$-step minimizes $J_h$ plus a quadratic penalty and
        decouples pointwise.
  \item The $(\bar\rho, b)$-step is a projection onto the affine constraint
        set $C$, which reduces to solving a single linear system.
\end{itemize}

%============================================================
\section{Linearized ADMM}
%============================================================

\subsection{Why Linearization is Necessary}

Standard (non-linearized) ADMM requires solving
\begin{equation}
  (\rho^{k+1}, m^{k+1})
  = \arg\min_{\rho,m}
    J_h(\rho,m)
    + \frac{\gamma}{2}\|A(\rho,m) - (\bar\rho^k, b^k) + \delta^k/\gamma\|^2.
  \label{eq:standard_admm}
\end{equation}
Because $A$ mixes neighboring values of $\rho$ in time and neighboring values
of $m$ in space, this is a \emph{coupled} minimization over the entire grid;
the quadratic term in~\eqref{eq:standard_admm} cannot be solved pointwise.

\subsection{The Linearization}

Linearized ADMM replaces the $A(\rho,m)$ term inside the penalty with a
first-order Taylor expansion around the current iterate, adding a proximal
term with parameter $\tau > 0$:
\begin{equation}
  (\rho^{k+1}, m^{k+1})
  = \operatorname{prox}_{J_h/\tau}\!\left(
      (\rho^k, m^k) - \frac{\gamma}{\tau}A^\top
      \!\left(A(\rho^k,m^k) - (\bar\rho^k, b^k) + \frac{\delta^k}{\gamma}
      \right)
    \right).
  \label{eq:linearized_step1}
\end{equation}
The argument to the proximal operator is an explicit function of the current
iterate, so no coupled solve is needed.  The proximal operator itself
decouples pointwise (see Section~\ref{sec:prox}).

\begin{remark}
  The stability condition for linearized ADMM is $\tau > \gamma\|A\|^2$.
  Since $A_\rho$ and $A_m$ are averaging operators, $\|A\|^2 \leq 1$, so
  $\tau = 5\gamma$ is a safe, practical choice.
\end{remark}

\subsection{The Three ADMM Steps}

Let $\delta = (\delta_{\bar\rho}, \delta_b)$ be the dual variable (Lagrange
multiplier for the ADMM constraint).

\paragraph{Step 1 --- $(\rho,m)$ proximal update.}
Compute the residual
$r = A(\rho^k,m^k) - (\bar\rho^k, b^k) + \delta^k/\gamma$,
form the gradient step
\begin{equation}
  u = \rho^k - \frac{\gamma}{\tau}A_\rho^\top r_{\bar\rho}, \qquad
  v = m^k   - \frac{\gamma}{\tau}A_m^\top r_b,
\end{equation}
and then compute $(\rho^{k+1}, m^{k+1}) = \operatorname{prox}_{J_h/\tau}(u, v)$
pointwise (Section~\ref{sec:prox}).

\paragraph{Step 2 --- $(\bar\rho, b)$ projection update.}
With over-relaxation parameter $\alpha \in (1,2)$:
\begin{equation}
  (\bar\rho^{k+1}, b^{k+1})
  = \operatorname{proj}_C\!\left(
      \alpha\,A(\rho^{k+1},m^{k+1}) + (1-\alpha)(\bar\rho^k,b^k)
      + \delta^k/\gamma
    \right).
\end{equation}
The projection solves a structured linear system (Section~\ref{sec:projection}).

\paragraph{Step 3 --- Dual update.}
\begin{equation}
  \delta^{k+1}
  = \delta^k + \gamma\bigl(
      \alpha\,A(\rho^{k+1},m^{k+1}) + (1-\alpha)(\bar\rho^k,b^k)
      - (\bar\rho^{k+1}, b^{k+1})
    \bigr).
\end{equation}

%============================================================
\section{Pointwise Proximal Operator}\label{sec:prox}
%============================================================

For each cell $(i,j)$, the proximal sub-problem is
\begin{equation}
  (\rho^*, m^*)
  = \arg\min_{\rho \geq 0,\, m}
    \frac{\Delta t\,\Delta x}{2\tau}\,\frac{m^2}{\rho}
    + \frac{1}{2}(\rho - u)^2 + \frac{1}{2}(m - v)^2.
  \label{eq:prox_problem}
\end{equation}
Setting $w = \Delta t\,\Delta x$, optimality in $m$ gives
\begin{equation}
  m^* = \frac{\tau\,\rho^*\,v}{\tau\,\rho^* + w/2}.
  \label{eq:mstar}
\end{equation}
Substituting back into~\eqref{eq:prox_problem} and differentiating with
respect to $\rho$ yields a \emph{cubic equation}:
\begin{equation}
  \tau^2 \rho^3
  - \tau^2 u\,\rho^2
  + \tau\!\left(w - \frac{w}{2}\right)\!(2w - u\,\tau)\,\rho
  - uw^2/4
  - \tfrac{w\tau}{2}\,v^2/4 = 0.
  \label{eq:cubic}
\end{equation}
More precisely, the coefficients are
\begin{equation}
  a_3 = \tau^2, \quad
  a_2 = \tau(2w - u\tau), \quad
  a_1 = w(w - 2u\tau), \quad
  a_0 = -uw^2 - \tfrac{w\tau}{2}\,v^2.
\end{equation}
We take the \textbf{largest real root} $\rho^* > 0$ (following Papadakis
et al.\ 2014).  If $u \leq 0$ we set $\rho^* = 0$, $m^* = 0$, since the
domain of $J_h$ requires $\rho > 0$.  A small positivity floor
$\rho^* \leftarrow \max(\rho^*, 10^{-10})$ is applied after the solve to
prevent division by zero in subsequent iterations.

The cubic is solved analytically (Cardano's method) in fully vectorized
NumPy code, so there are no Python loops over grid points.

%============================================================
\section{Implicit Diffusion and the CFL Constraint}\label{sec:implicit}
%============================================================

\subsection{The Parabolic CFL Condition}

An \emph{explicit} (forward-Euler) discretization of the diffusion term
$\varepsilon\,\partial_{xx}\rho$ would give
\[
  \frac{\bar\rho^{n+1} - \bar\rho^n}{\Delta t}
  = \frac{\varepsilon}{\Delta x^2}(\bar\rho^n_{j+1} - 2\bar\rho^n_j + \bar\rho^n_{j-1}).
\]
Von Neumann stability analysis requires the \emph{parabolic Courant number}
$\mu = \varepsilon\Delta t/\Delta x^2 \leq 1/2$.  For fixed spatial
resolution $N_x$ this imposes
\begin{equation}
  \Delta t \leq \frac{\Delta x^2}{2\varepsilon}
  = \frac{1}{2\varepsilon N_x^2}.
  \label{eq:CFL}
\end{equation}
For the parameters $N_x = 128$, $\varepsilon = 0.05$ this gives
$\Delta t \leq 6\times10^{-4}$, requiring $N_t \geq 1670$ time steps.

\subsection{Our Implicit Treatment}

The Fokker--Planck constraint~\eqref{eq:FP} is discretized as
\begin{equation}
  (I_x \otimes D_t - \varepsilon L_x \otimes A_t)\,\bar\rho
  + (D_x \otimes I_t)\,b = d,
  \label{eq:discrete_FP}
\end{equation}
where:
\begin{itemize}
  \item $D_t$ (shape $N_t \times (N_t-1)$) is the first-difference operator
        in time, corresponding to the time derivative $\partial_t\rho$.
  \item $A_t$ (shape $N_t \times (N_t-1)$) is the time-averaging operator.
        The diffusion term is evaluated at the \emph{average} of the
        density at adjacent integer-time nodes: this is an
        \textbf{implicit/semi-implicit (Crank-Nicolson-like) scheme}.
  \item $L_x = D_x D_x^\top$ is the negative Neumann Laplacian
        ($L_x \approx -\partial_{xx}$, positive semi-definite), so
        $+\varepsilon L_x$ correctly represents forward diffusion.
  \item $D_x \otimes I_t$ is the spatial divergence of the flux $b$.
\end{itemize}

Because $A_t$ averages between two adjacent time levels, the diffusion term
involves $\bar\rho$ at both $t_i$ and $t_{i+1}$.  This is equivalent to an
implicit treatment of the diffusion, making the scheme
\textbf{unconditionally stable} with respect to the diffusion coefficient.
There is no restriction of the form~\eqref{eq:CFL}.

\subsection{Stiffness Ratio}

A useful diagnostic is the \emph{stiffness ratio}
\begin{equation}
  \mu_h = \frac{\varepsilon}{\Delta x^2 / \Delta t}
        = \frac{\varepsilon N_x^2}{N_t}.
  \label{eq:stiffness}
\end{equation}
When $\mu_h \gg 1$ the scheme is operating far outside the explicit stability
region, yet convergence is maintained because the implicit diffusion enters
the Fokker--Planck constraint solve (Step 2 of ADMM) rather than the time
integration.  In practice, values of $\mu_h$ up to $\sim 10$ are handled
without difficulty.

\begin{remark}
  The advective CFL $|v|\Delta t/\Delta x \leq 1$ does not appear explicitly
  because the velocity $v = m/\rho$ is not discretized directly.  Instead,
  $m$ is the independent variable and the kinetic energy $m^2/(2\rho)$ is
  minimized directly.
\end{remark}

%============================================================
\section{The Projection Step}\label{sec:projection}
%============================================================

\subsection{KKT System}

The projection
$(\bar\rho^{k+1}, b^{k+1}) = \operatorname{proj}_C(p_{\bar\rho}, p_b)$
is the nearest point in $C$ to $(p_{\bar\rho}, p_b)$.  By KKT conditions,
the solution satisfies
\begin{align}
  \bar\rho^{k+1} &= p_{\bar\rho}
    - (I_x \otimes D_t - \varepsilon L_x \otimes A_t)^\top \phi,
    \label{eq:KKT_rhobar}\\
  b^{k+1} &= p_b - (D_x \otimes I_t)^\top \phi
  \label{eq:KKT_b}
\end{align}
where the multiplier $\phi \in \mathbb{R}^{N_t \times N_x}$ solves the
normal equations
\begin{equation}
  M_{\mathrm{SB}}\,\phi = F,
  \label{eq:MSB_system}
\end{equation}
with right-hand side
$F = (I_x \otimes D_t - \varepsilon L_x \otimes A_t)\,p_{\bar\rho}
    + (D_x \otimes I_t)\,p_b - d$,
and the symmetric positive semi-definite operator
\begin{align}
  M_{\mathrm{SB}}
  &= (I_x \otimes D_t - \varepsilon L_x \otimes A_t)
     (I_x \otimes D_t - \varepsilon L_x \otimes A_t)^\top
   + (D_x \otimes I_t)(D_x \otimes I_t)^\top \nonumber\\
  &= I_x \otimes D_t D_t^\top
   - \varepsilon L_x \otimes (A_t D_t^\top + D_t A_t^\top)
   + \varepsilon^2 L_x^2 \otimes A_t A_t^\top
   + D_x D_x^\top \otimes I_t.
  \label{eq:MSB}
\end{align}
A direct assembly and factorization of $M_{\mathrm{SB}}$ would cost
$O(N_t^3 N_x^3)$.  We exploit its tensor-product structure for an
$O(N_x N_t \log N_t)$ solve.

\subsection{Efficient Solver via DCT and Tridiagonal Systems}

The key observation is that $L_x = D_x D_x^\top$ is \emph{diagonalized by
the DCT-II basis} (Neumann boundary conditions).  Applying DCT-II in space
decouples~\eqref{eq:MSB_system} into $N_x$ independent $N_t \times N_t$
tridiagonal systems, one per spatial mode $k = 0, \ldots, N_x - 1$:
\begin{equation}
  T^k \hat\phi_k = \hat F_k,
  \label{eq:tridiag}
\end{equation}
where $\hat F_k$ is the $k$-th DCT-II mode of $F$ and
$T^k$ is a symmetric tridiagonal matrix with entries
\begin{align}
  \text{main diagonal (interior):} &\quad
    \alpha = \frac{2}{\Delta t^2} + \frac{\varepsilon^2 \lambda_k^2}{2} + \lambda_k,\\
  \text{off-diagonal:} &\quad
    \beta  = -\frac{1}{\Delta t^2} + \frac{\varepsilon^2 \lambda_k^2}{4},\\
  \text{corner $(0,0)$:} &\quad
    \frac{1}{\Delta t^2} + \frac{\varepsilon\lambda_k}{\Delta t}
    + \frac{\varepsilon^2\lambda_k^2}{4} + \lambda_k,\\
  \text{corner $(N_t{-}1, N_t{-}1)$:} &\quad
    \frac{1}{\Delta t^2} - \frac{\varepsilon\lambda_k}{\Delta t}
    + \frac{\varepsilon^2\lambda_k^2}{4} + \lambda_k,
\end{align}
where $\lambda_k = \frac{2}{\Delta x^2}(1 - \cos(\pi k/N_x)) \geq 0$ is
the $k$-th eigenvalue of $L_x$.

For $k > 0$ (where $\lambda_k > 0$), $T^k$ is strictly positive definite
and is solved with SciPy's $O(N_t)$ banded solver.  For $k = 0$
(the zero-mode of $L_x$), $T^0 = D_t D_t^\top$ is the Neumann Laplacian
in time, which is singular; the right-hand side is guaranteed to be
consistent by mass conservation, and the minimum-norm (mean-zero) solution
is returned.

\paragraph{Computational cost.}
The DCT-II transform costs $O(N_x N_t \log N_x)$.  Each of the $N_x$
tridiagonal solves costs $O(N_t)$.  The inverse DCT-II costs
$O(N_x N_t \log N_x)$.  The total is
$O(N_x N_t (\log N_x + 1))$, which is essentially linear in the problem
size $N_x N_t$.

%============================================================
\section{Discrete Operators in Detail}
%============================================================

\subsection{Notation and Ordering}

All 2D arrays use the convention: \textbf{time is the slow index} (row
index) and \textbf{space is the fast index} (column index).  Kronecker
products are written with the time factor first:
$B_t \otimes B_x$ acts on a vector with time-slow ordering.

\subsection{One-Dimensional Operators}

Table~\ref{tab:operators} summarizes the 1D building blocks.

\begin{table}[h]
\centering
\caption{One-dimensional discrete operators.}
\label{tab:operators}
\begin{tabular}{llll}
\toprule
Symbol & Shape & Action & Notes \\
\midrule
$D_t$ & $N_t \times (N_t-1)$ & Time first-difference & BCs in $d$ \\
$A_t$ & $N_t \times (N_t-1)$ & Time averaging & Semi-implicit diffusion \\
$D_x$ & $N_x \times (N_x-1)$ & Space first-difference & Interior faces only \\
$L_x = D_x D_x^\top$ & $N_x \times N_x$ & Neumann Laplacian & $L_x \approx -\partial_{xx}$ \\
$A_\rho$ & $(N_t-1) \times N_t$ & Time avg: Grid~1$\to$Grid~2 & \eqref{eq:Arho} \\
$A_m$   & $(N_x-1) \times N_x$ & Space avg: Grid~1$\to$Grid~2 & \eqref{eq:Am} \\
\bottomrule
\end{tabular}
\end{table}

For $N_t = 4$ the time operators are (after scaling by $\Delta t$):
\[
  D_t = \frac{1}{\Delta t}\begin{pmatrix}
    1 &  0 &  0 \\
   -1 &  1 &  0 \\
    0 & -1 &  1 \\
    0 &  0 & -1
  \end{pmatrix},
  \qquad
  A_t = \frac{1}{2}\begin{pmatrix}
    1 & 0 & 0 \\
    1 & 1 & 0 \\
    0 & 1 & 1 \\
    0 & 0 & 1
  \end{pmatrix}.
\]

\subsection{Right-Hand Side Vector $d$}

The boundary conditions $\rho(0,\cdot) = \rho_0$, $\rho(1,\cdot) = \rho_1$
are absorbed into the right-hand side:
\begin{equation}
  d = \frac{1}{\Delta t}
  \begin{pmatrix} \rho_0 \\ 0 \\ \vdots \\ 0 \\ -\rho_1 \end{pmatrix},
  \label{eq:rhs}
\end{equation}
with $d \in \mathbb{R}^{N_t \times N_x}$ (time-slow ordering).  The
factor $1/\Delta t$ comes from the first-difference operator $D_t$ acting
on a constant vector.

%============================================================
\section{Algorithm Summary}
%============================================================

\begin{algorithm}[h]
\caption{Two-grid Linearized ADMM for the Schr\"{o}dinger Bridge}
\label{alg:admm}
\begin{algorithmic}[1]
\Require $\rho_0$, $\rho_1$: boundary densities; $N_t$, $N_x$, $\varepsilon$, $\gamma$, $\tau$, $\alpha$
\State Build operators $D_t$, $A_t$, $D_x$, $L_x$, $A_\rho$, $A_m$ and RHS $d$
\State Initialize $\rho$ by center-of-mass displacement interpolation of $\rho_0$
\State $m \leftarrow (\mu_1-\mu_0)\,\rho$ \quad (constant-velocity momentum)
\State $\bar\rho \leftarrow A_\rho\,\rho$, \quad $b \leftarrow A_m\,m$, \quad $\delta_{\bar\rho} \leftarrow 0$, \quad $\delta_b \leftarrow 0$
\For{$k = 1, 2, \ldots$}
  \State \textbf{Step 1: Proximal update}
  \State $r_{\bar\rho} \leftarrow A_\rho\rho - \bar\rho + \delta_{\bar\rho}/\gamma$, \quad $r_b \leftarrow A_m m - b + \delta_b/\gamma$
  \State $u \leftarrow \rho - (\gamma/\tau)\,A_\rho^\top r_{\bar\rho}$, \quad $v \leftarrow m - (\gamma/\tau)\,A_m^\top r_b$
  \State $(\rho, m) \leftarrow \operatorname{prox}_{J_h/\tau}(u, v)$ \quad (pointwise cubic solve, \S\ref{sec:prox})
  \State \textbf{Step 2: Projection update (with over-relaxation)}
  \State $p_{\bar\rho} \leftarrow \alpha A_\rho\rho + (1-\alpha)\bar\rho + \delta_{\bar\rho}/\gamma$
  \State $p_b \leftarrow \alpha A_m m + (1-\alpha) b + \delta_b/\gamma$
  \State $(\bar\rho, b) \leftarrow \operatorname{proj}_C(p_{\bar\rho}, p_b)$ \quad (DCT + tridiagonal solve, \S\ref{sec:projection})
  \State \textbf{Step 3: Dual update}
  \State $\delta_{\bar\rho} \leftarrow \delta_{\bar\rho} + \gamma(p_{\bar\rho} - \bar\rho)$, \quad $\delta_b \leftarrow \delta_b + \gamma(p_b - b)$
  \State Check primal/dual residuals; break if $< \texttt{tol}$
\EndFor
\end{algorithmic}
\end{algorithm}

%============================================================
\section{Convergence Diagnostics}
%============================================================

At each iteration we track:
\begin{enumerate}
  \item \textbf{Primal residual} (normalized):
    $r_p = \|A(\rho^k,m^k) - (\bar\rho^k,b^k)\| / \sqrt{N_{\mathrm{primal}}}$
  \item \textbf{Dual residual} (normalized):
    $r_d = \gamma\,\|A^\top((\bar\rho^k,b^k) - (\bar\rho^{k-1},b^{k-1}))\| / \sqrt{N_{\mathrm{dual}}}$
  \item \textbf{Objective value}:
    $J_h(\rho^k, m^k) = \Delta t\,\Delta x\sum_{ij} m_{ij}^2/(2\rho_{ij})$
  \item \textbf{Fokker--Planck violation}:
    $\|(I_x \otimes D_t - \varepsilon L_x \otimes A_t)\bar\rho^k
       + (D_x \otimes I_t)b^k - d\|$
\end{enumerate}
Convergence is declared when $r_p < \texttt{tol}$ (default $10^{-4}$).
The FP violation is an independent check that the projection is working
correctly; it should be near machine precision at every iteration.

%============================================================
\section{Connection to the Literature}
%============================================================

This formulation follows the line of work by Benamou and Brenier (2000)
who introduced the dynamic OT problem~\eqref{eq:BB_objective}--\eqref{eq:continuity},
and the ADMM discretization approach of Papadakis, Peyr\'{e}, and Oudet (2014)
who developed the proximal splitting on two grids for the OT case ($\varepsilon=0$).

The extension to $\varepsilon > 0$ (Schr\"{o}dinger Bridge) adds the
diffusion term $\varepsilon L_x \otimes A_t$ to the Fokker--Planck constraint
and modifies the $M_{\mathrm{SB}}$ operator accordingly.  The semi-implicit
treatment of this term via $A_t$ (averaging in time) is the key stability
contribution: it is the discrete analogue of the Crank--Nicolson scheme for
the heat equation but embedded inside the ADMM projection step rather than
an explicit time-stepper.

For Gaussian boundary conditions, the exact SB path has an analytic formula:
$\rho(t,\cdot) = \mathcal{N}(\mu(t), \sigma^2(t))$ where $\mu(t)$ is linear
in $t$ and $\sigma^2(t)$ follows from the cross-covariance
$c = -\varepsilon T + \sqrt{(\varepsilon T)^2 + \sigma_0^2\sigma_1^2}$.
This provides a ground-truth comparison for the numerical solver.

%============================================================
\section{Parameters and Test Cases}
%============================================================

\subsection{Recommended Parameters}

\begin{table}[h]
\centering
\caption{Recommended parameter settings.}
\begin{tabular}{lll}
\toprule
Parameter & Default & Notes \\
\midrule
$N_t$ & 64 & Time cells \\
$N_x$ & 128 & Space cells \\
$\varepsilon$ & $0$--$0.05$ & $0$ = pure OT \\
$\gamma$ & 1.0 & ADMM penalty \\
$\tau$   & 5.0 & Proximal parameter; need $\tau > \gamma\|A\|^2 \approx \gamma$ \\
$\alpha$ & 1.5 & Over-relaxation (Boyd et al.\ 2011) \\
Tolerance & $10^{-4}$ & On normalized primal residual \\
\bottomrule
\end{tabular}
\end{table}

\subsection{Test Cases}

\paragraph{Test 1 --- Gaussian $\to$ Gaussian, $\varepsilon = 0$ (OT):}
$\rho_0 = \mathcal{N}(0.25, 0.05^2)$, $\rho_1 = \mathcal{N}(0.75, 0.05^2)$.
Expected objective: $J^* = \frac{1}{2}(0.75 - 0.25)^2 = 0.125$.
The density evolves as a linear displacement interpolation.

\paragraph{Test 2 --- Same marginals, $\varepsilon = 0.05$ (SB):}
Expect smoother, wider paths; $J_{\mathrm{SB}} > J_{\mathrm{OT}}$ due to
the entropic cost.  Comparison against the analytic Gaussian formula
validates the solver.

\paragraph{Test 3 --- Two humps $\to$ one hump:}
$\rho_0 = \frac{1}{2}\mathcal{N}(0.25,0.05^2) + \frac{1}{2}\mathcal{N}(0.75,0.05^2)$,
$\rho_1 = \mathcal{N}(0.5, 0.1^2)$.
Tests behavior with bimodal initial data.

%============================================================
\appendix
\section{Why Not Discretize the HJB Equation?}
%============================================================

An alternative approach is to derive the optimality conditions of the
continuous problem first.  The Pontryagin maximum principle yields a coupled
system of a forward Fokker--Planck equation for $\rho$ and a backward
Hamilton--Jacobi equation for the value function $\phi$, linked through the
optimal velocity $v^* = \partial_x \phi$.  This \emph{Schr\"{o}dinger system}
can be solved iteratively (e.g., Sinkhorn-like fixed-point iterations in the
static setting).

The DtO approach chosen here has several practical advantages over this route:
\begin{enumerate}
  \item \textbf{No adjoint PDE.}  The ADMM objective encodes the constraints
        implicitly; there is no backward PDE to discretize consistently.
  \item \textbf{Arbitrary boundary conditions.}  Periodic, Neumann, and
        Dirichlet boundary conditions are handled identically by replacing
        the operators $D_t$, $D_x$, $L_x$ without changing the algorithm.
  \item \textbf{Direct access to $(\rho, m)$.}  The primal variables are
        the physically meaningful density and momentum, not potentials.
        Positivity of $\rho$ is enforced at every iteration via the proximal
        operator.
  \item \textbf{No CFL constraint.}  As explained in Section~\ref{sec:implicit},
        the implicit diffusion in the projection step sidesteps the parabolic
        stability restriction entirely.
\end{enumerate}

%============================================================
\begin{thebibliography}{9}
%============================================================

\bibitem{BB2000}
  J.-D.\ Benamou and Y.\ Brenier,
  ``A computational fluid mechanics solution to the Monge-Kantorovich
  mass transfer problem,''
  \textit{Numerische Mathematik}, 84(3):375--393, 2000.

\bibitem{Papadakis2014}
  N.\ Papadakis, G.\ Peyr\'{e}, and E.\ Oudet,
  ``Optimal transport with proximal splitting,''
  \textit{SIAM Journal on Imaging Sciences}, 7(1):212--238, 2014.

\bibitem{Boyd2011}
  S.\ Boyd, N.\ Parikh, E.\ Chu, B.\ Peleato, and J.\ Eckstein,
  ``Distributed optimization and statistical learning via the alternating
  direction method of multipliers,''
  \textit{Foundations and Trends in Machine Learning}, 3(1):1--122, 2011.

\bibitem{Leonard2014}
  C.\ L\'{e}onard,
  ``A survey of the Schr\"{o}dinger problem and some of its connections
  with optimal transport,''
  \textit{Discrete and Continuous Dynamical Systems}, 34(4):1533--1574, 2014.

\end{thebibliography}

\end{document}
